\chapter{Diagramas de Cableado}

Este capítulo muestra ejemplos básicos de cableado para proyectos comunes usando el firmware del Controlador RC con ESP32.

Estos diagramas están simplificados para que los principiantes puedan entender rápidamente cómo conectar el hardware.

\section{Cableado Básico de Robot}

El proyecto más común es un robot de tracción diferencial simple usando dos motores.

% Requiere: \usetikzlibrary{positioning,arrows.meta,calc}
% (Ya cargas positioning+arrows.meta en main.tex.)
% Diseño más angosto (menos ancho) y opcionalmente más alto.
% Usa un flujo mayormente vertical: Teléfono -> ESP32 -> Driver, con motores ramificando a la derecha.
% Requiere: \usetikzlibrary{positioning,arrows.meta,calc}

\begin{center}
    \begin{tikzpicture}[
        font=\small,
        >=Stealth,
        node distance=22mm and 34mm, % espaciado más ajustado (vertical y horizontal)
        box/.style={
            draw,
            rounded corners=2mm,
            very thick,
            align=center,
            minimum width=36mm,  % ancho reducido
            minimum height=14mm, % ligeramente más alto
            fill=gray!6
        },
        boxAccent/.style={
            box,
            fill=blue!6,
            draw=blue!60!black
        },
        boxHW/.style={
            box,
            fill=green!6,
            draw=green!50!black
        },
        arrow/.style={->, very thick, draw=black!70},
        label/.style={midway, fill=white, inner sep=1.5pt, text=black!80, font=\scriptsize}
        ]

        % Nodos (apilados verticalmente para reducir el ancho)
        \node[boxAccent] (phone) {App\\Android};
        \node[boxAccent, below=of phone] (esp32) {ESP32\\Controlador};
        \node[boxHW, below=of esp32] (driver) {Driver\\de Motor};

        % Los motores se ramifican a la derecha (compacto)
        \node[boxHW, right=16mm of driver, yshift=+9mm] (motorL) {Motor\\Izquierdo};
        \node[boxHW, right=16mm of driver, yshift=-9mm] (motorR) {Motor\\Derecho};

        % Flechas + etiquetas
        \draw[arrow] (phone) -- node[label, right]{Bluetooth SPP} (esp32);
        \draw[arrow] (esp32) -- node[label, right, yshift=-3mm]{PWM + DIR} (driver);

        \draw[arrow] (driver.east) -- node[label, above, sloped]{Potencia} (motorL.west);
        \draw[arrow] (driver.east) -- node[label, below, sloped]{Potencia} (motorR.west);

        % Agrupación opcional sutil
        \draw[dashed, rounded corners=2mm, draw=black!30]
        ($(phone.north west)+(-4mm,4mm)$) rectangle ($(esp32.south east)+(4mm,-4mm)$);
        \node[font=\normalsize, text=black!60] at ($(phone.west)!0.5!(esp32.north)-(3.5,6mm)$) {Enlace de Control};

        \draw[dashed, rounded corners=2mm, draw=black!30]
        ($(driver.north west)+(-4mm,12mm)$) rectangle ($(motorR.south east)+(4mm,-4mm)$);
        \node[font=\normalsize, text=black!60] at ($(driver.north)!0.55!(motorL.north)+(11mm,10mm)$) {Etapa de Potencia};

    \end{tikzpicture}
\end{center}

\section{Ejemplo de Control de Servo}

Los servos pueden conectarse directamente a pines PWM del ESP32.

\begin{itemize}
    \item Señal del servo → pin GPIO del ESP32
    \item Alimentación del servo → fuente de 5V
    \item Tierra del servo → tierra común
\end{itemize}

Ejemplos de uso:

\begin{itemize}
    \item Control de dirección
    \item Control de paneo de cámara
    \item Articulaciones de brazo robótico
\end{itemize}

\section{Salidas para LED y Zumbador}

Las salidas digitales pueden controlar LEDs o zumbadores.

Conexiones de ejemplo:

\begin{itemize}
    \item Pin GPIO → LED (con resistencia)
    \item Pin GPIO → módulo zumbador
\end{itemize}

Estas salidas pueden controlarse usando botones o interruptores en la aplicación Android.

\section{Conexiones de Sensores}

Los sensores pueden conectarse usando:

\begin{itemize}
    \item Entradas analógicas
    \item Entradas digitales
    \item Interfaz I2C
\end{itemize}

El firmware permite transmitir datos de sensores a la aplicación Android mediante telemetría.